\documentclass[11pt,a4paper]{article}

\usepackage[margin=1in]{geometry}
\usepackage{graphicx}
\usepackage{booktabs}
\usepackage{tabularx}
\usepackage{hyperref}
\usepackage{xcolor}
\usepackage{listings}
\usepackage{float}
\usepackage{enumitem}
\usepackage{fancyhdr}
\usepackage{titlesec}
\usepackage{setspace}
\usepackage[T1]{fontenc}
\usepackage{longtable}

% ── Color definitions ──
\definecolor{codeblue}{HTML}{3B82F6}
\definecolor{codegray}{HTML}{64748B}
\definecolor{codegreen}{HTML}{10B981}
\definecolor{codered}{HTML}{EF4444}
\definecolor{codebg}{HTML}{F1F5F9}
\definecolor{sectionblue}{HTML}{1E40AF}

\lstset{
  basicstyle=\ttfamily\small,
  backgroundcolor=\color{codebg},
  keywordstyle=\color{codeblue}\bfseries,
  commentstyle=\color{codegray}\itshape,
  stringstyle=\color{codegreen},
  breaklines=true,
  frame=single,
  rulecolor=\color{codegray!30},
  numbers=left,
  numberstyle=\tiny\color{codegray},
  numbersep=8pt,
  tabsize=4,
  showstringspaces=false,
  captionpos=b,
}

\pagestyle{fancy}
\fancyhf{}
\fancyhead[L]{\textsc{smoltcp Master Test Plan}}
\fancyhead[R]{\thepage}
\fancyfoot[C]{\small\textit{Test Planning Document}}

\titleformat{\section}{\Large\bfseries\color{sectionblue}}{\thesection}{1em}{}
\titleformat{\subsection}{\large\bfseries\color{sectionblue!80}}{\thesubsection}{1em}{}
\titleformat{\subsubsection}{\normalsize\bfseries\color{sectionblue!60}}{\thesubsubsection}{1em}{}

\setstretch{1.15}

\begin{document}

% ══════════════════════════════════════════════════════════════════
\begin{titlepage}
\centering
\vspace*{2cm}
{\Huge\bfseries\color{sectionblue} Master Test Plan\\[0.4cm] smoltcp v0.12.0\par}
\vspace{1.5cm}
{\Large Security \& Quality Testing Campaign\par}
\vspace{1cm}
{\large TCP/IP Stack Quality Assurance\par}
\vspace{2cm}
{\large\textit{Test Planning Document}\par}
\vspace{0.5cm}
{\large April 2026\par}
\vfill
{\small Document version: 1.0\par}
{\small Platforms: Windows 10 (MSVC) | WSL Ubuntu 24.04.2 LTS\par}
\end{titlepage}

\tableofcontents
\newpage

% ══════════════════════════════════════════════════════════════════
\section{Test Plan Overview}
% ══════════════════════════════════════════════════════════════════

\subsection{Purpose}

This Master Test Plan outlines the strategy, environment, resources, schedule, and evaluation criteria for a comprehensive quality assessment of smoltcp v0.12.0, an open-source, event-driven TCP/IP stack written in Rust. The plan coordinates eight distinct testing phases spanning mutation adequacy, functional testing, security analysis, conformance testing, performance evaluation, white-box adequacy, safety assessment, and defect correction.

\subsection{Scope}

Testing focuses on four high-complexity source modules selected for their critical role in packet processing:

\begin{itemize}[noitemsep]
\item \texttt{src/storage/assembler.rs} --- TCP segment reassembly with gap tracking
\item \texttt{src/wire/ipv4.rs} --- IPv4 header parsing, validation, and emission
\item \texttt{src/wire/udp.rs} --- UDP datagram parsing and checksum verification
\item \texttt{src/wire/tcp.rs} --- TCP segment parsing, flag extraction, and option handling
\end{itemize}

\subsection{Standards and References}

\begin{itemize}[noitemsep]
\item RFC 791 --- Internet Protocol (IPv4)
\item RFC 768 --- User Datagram Protocol (UDP)
\item RFC 793 --- Transmission Control Protocol (TCP)
\item RFC 2460 --- Internet Protocol, Version 6 (IPv6)
\item IEEE 802.15.4 --- Low-Rate Wireless Personal Area Networks
\item IETF RFC 6282 --- 6LoWPAN Header Compression
\end{itemize}

% ══════════════════════════════════════════════════════════════════
\section{Test Objectives}
% ══════════════════════════════════════════════════════════════════

The overall testing objectives are:

\begin{enumerate}[noitemsep]
\item \textbf{Evaluate test adequacy} via mutation analysis on four complex modules and quantify kill ratios per module and per operator class.
\item \textbf{Improve test suite} based on surviving non-equivalent mutants by adding targeted assertions.
\item \textbf{Perform functional testing} using the category-partition method with combinatorial test generation for UDP and IPv4 parsing.
\item \textbf{Conduct security fuzzing} with coverage-guided mutation (libFuzzer) to identify crash-inducing inputs.
\item \textbf{Verify conformance} against RFC 791 and RFC 768 header parsing requirements.
\item \textbf{Benchmark performance} using loopback throughput measurement and microbenchmarks.
\item \textbf{Measure white-box coverage} using LLVM source-based coverage instrumentation.
\item \textbf{Assess software safety} through panic-safety testing and unsafe code inventory.
\item \textbf{Apply defect corrections} for any discovered specification deviations.

\end{enumerate}

% ══════════════════════════════════════════════════════════════════
\section{Test Environment}
% ══════════════════════════════════════════════════════════════════

\subsection{Hardware}

All testing is conducted on a single host machine with the following specifications:
\begin{itemize}[noitemsep]
\item CPU: AMD/Intel x86\_64 processor
\item RAM: 16 GB minimum
\item Storage: SSD with $\geq$50 GB free space for build artifacts
\end{itemize}

\subsection{Software Configurations}

Two platform configurations are tested:

\begin{table}[H]
\centering
\caption{Test environment configurations.}
\begin{tabular}{lll}
\toprule
\textbf{Attribute} & \textbf{Windows} & \textbf{WSL} \\
\midrule
Operating System   & Windows 10 (64-bit) & Ubuntu 24.04.2 LTS \\
Shell              & PowerShell 7.x      & Bash 5.x \\
Rust Stable        & 1.91+               & 1.91+ \\
Rust Nightly       & 1.96.0-nightly      & 1.96.0-nightly \\
C Compiler         & MSVC (cl.exe)       & GCC/Clang \\
cargo-fuzz         & 0.13.1              & 0.13.1 \\
cargo-llvm-cov     & 0.8.5               & 0.8.5 \\
Python             & 3.12+               & 3.12+ \\
\bottomrule
\end{tabular}
\end{table}

\subsection{Configuration Strategy}

Given resource constraints, the following testing strategy is applied:
\begin{itemize}[noitemsep]
\item Both platforms are tested extensively (``extreme point'' approach).
\item Core functional tests (mutation, partitioning, conformance, safety, coverage) run on both platforms.
\item Fuzzing is Linux-only due to MSVC/libFuzzer incompatibility.
\item Performance results are compared across platforms to identify OS-specific bottlenecks.
\end{itemize}

% ══════════════════════════════════════════════════════════════════
\section{Test Suites}
% ══════════════════════════════════════════════════════════════════

\subsection{Suite 1: Mutation Adequacy}

\begin{itemize}[noitemsep]
\item \textbf{Objective:} Evaluate test suite effectiveness via first-order mutation analysis.
\item \textbf{Method:} Custom Python mutant generator applying 7 operator classes.
\item \textbf{Target files:} 4 source modules, $\geq$300 total mutants.
\item \textbf{Oracle:} Module-level unit test suite (cargo test --lib).
\item \textbf{Metrics:} Kill ratio $K/(M-E)$ per file and overall.
\item \textbf{Deliverables:} \texttt{mutants.diff}, \texttt{mutation\_results\_\{platform\}.csv}, per-mutant execution logs.
\end{itemize}

\begin{lstlisting}[language=bash, caption={Mutation testing invocation.}]
python mutate_and_test.py --source src/storage/assembler.rs \
    --operators rel,arith,bool,const,bit,assign,shift \
    --output-dir logs/ --platform windows
\end{lstlisting}

\subsection{Suite 2: Input Space Partitioning (UDP)}

\begin{itemize}[noitemsep]
\item \textbf{Objective:} Systematic functional testing of UDP parsing via the category-partition method.
\item \textbf{Method:} Full combinatorial coverage of 6 input dimensions.
\item \textbf{Total cases:} 480 test frames.
\item \textbf{Oracle:} Expected \texttt{new\_checked()} and \texttt{parse()} outcomes derived from RFC 768 and smoltcp behavior.
\item \textbf{Deliverables:} \texttt{udp\_cases.csv}, partition suite source code, execution logs.
\end{itemize}

\subsection{Suite 3: IPv4 Combinatorial}

\begin{itemize}[noitemsep]
\item \textbf{Objective:} Category-partition testing of IPv4 header parsing.
\item \textbf{Method:} Combinatorial coverage of version, IHL, total length, fragmentation, and checksum.
\item \textbf{Total cases:} 216 test frames.
\item \textbf{Oracle:} Aligned to smoltcp's actual fragmentation and header-length behavior.
\item \textbf{Deliverables:} \texttt{ipv4\_cases\_\{platform\}.csv}, JSON summary.
\end{itemize}

\subsection{Suite 4: Security Fuzzing}

\begin{itemize}[noitemsep]
\item \textbf{Objective:} Discover crash-inducing inputs via coverage-guided mutation fuzzing.
\item \textbf{Method:} libFuzzer (via cargo-fuzz) with sanitizer instrumentation.
\item \textbf{Targets:} 5 fuzz harnesses (\texttt{packet\_parser}, \texttt{tcp\_headers}, \texttt{dhcp\_header}, \texttt{ieee802154\_header}, \texttt{sixlowpan\_packet}).
\item \textbf{Time budget:} 60s initial + 180s extended per target.
\item \textbf{Oracle:} Process exit code (crash = failure); sanitizer reports for root cause.
\item \textbf{Deliverables:} Crash artifacts, reproduction logs, coverage statistics.
\end{itemize}

\subsection{Suite 5: Conformance Testing}

\begin{itemize}[noitemsep]
\item \textbf{Objective:} Verify adherence to RFC 791 (IPv4) and RFC 768 (UDP).
\item \textbf{Method:} Crafted packets testing specific header validation rules.
\item \textbf{Total cases:} 9 hand-crafted conformance assertions.
\item \textbf{Oracle:} RFC-derived expected accept/reject decisions.
\item \textbf{Deliverables:} \texttt{conformance\_results.csv}, per-platform logs.
\end{itemize}

\subsection{Suite 6: Performance Testing}

\begin{itemize}[noitemsep]
\item \textbf{Objective:} Measure raw throughput and per-operation latency.
\item \textbf{Methods:} Loopback throughput benchmark (128 MB transfer, 5 runs) and microbenchmarks (5 repetitions).
\item \textbf{Metrics:} Gbps throughput, ns/iteration latency.
\item \textbf{Deliverables:} CSV metrics per platform, execution logs.
\end{itemize}

\subsection{Suite 7: White-Box Adequacy}

\begin{itemize}[noitemsep]
\item \textbf{Objective:} Quantify source code coverage.
\item \textbf{Tool:} \texttt{cargo-llvm-cov v0.8.5} (LLVM source-based coverage).
\item \textbf{Metrics:} Region, function, and line coverage percentages.
\item \textbf{Deliverables:} Coverage reports, LCOV artifacts, per-module breakdowns.
\end{itemize}

\subsection{Suite 8: Software Safety}

\begin{itemize}[noitemsep]
\item \textbf{Objective:} Verify panic-free operation on invalid inputs and audit unsafe code usage.
\item \textbf{Methods:} Panic-safety test suite; \texttt{grep -r ``unsafe''} inventory with manual classification.
\item \textbf{Deliverables:} Safety test logs, \texttt{unsafe\_inventory.csv}.
\end{itemize}

% ══════════════════════════════════════════════════════════════════
\section{Evaluation Criteria}
% ══════════════════════════════════════════════════════════════════

\begin{table}[H]
\centering
\caption{Evaluation criteria for each test dimension.}
\begin{tabularx}{\textwidth}{lX}
\toprule
\textbf{Dimension} & \textbf{Criteria} \\
\midrule
Mutation adequacy     & Kill ratio $K/(M-E) \geq 70\%$ per file; overall $\geq 65\%$ \\
Input partitioning    & Pairwise coverage achieved; $\geq 95\%$ tests passing \\
Fuzzing               & $\geq 6$ distinct error signatures documented (or justification if fewer) \\
Conformance           & All tests pass or deviations documented with rationale \\
Performance           & Median throughput and variance reported; cross-platform comparison \\
White-box adequacy    & $\geq 80\%$ line coverage for target modules \\
Safety                & No panics on invalid inputs; all unsafe usage reviewed \\
\bottomrule
\end{tabularx}
\end{table}

% ══════════════════════════════════════════════════════════════════
\section{Test Execution Schedule}
% ══════════════════════════════════════════════════════════════════

Testing follows a sequential single-threaded execution model:

\begin{enumerate}[noitemsep]
\item \textbf{Phase 1 --- Infrastructure:} Create test suites, harnesses, and automation scripts.
\item \textbf{Phase 2 --- Windows execution:} Run all suites sequentially on Windows (PowerShell).
\item \textbf{Phase 3 --- WSL execution:} Run all suites sequentially on WSL (Bash).
\item \textbf{Phase 4 --- Enhanced suites:} Execute IPv4 combinatorial, microbenchmarks, enhanced fuzzing, and white-box additions.
\item \textbf{Phase 5 --- Analysis:} Analyze surviving mutants, compare cross-platform results, identify defects.
\item \textbf{Phase 6 --- Remediation:} Apply defect corrections, re-run relevant suites to verify fixes.
\item \textbf{Phase 7 --- Reporting:} Generate final report, test plan, and presentation.
\end{enumerate}

% ══════════════════════════════════════════════════════════════════
\section{Test Deliverables}
% ══════════════════════════════════════════════════════════════════

\begin{itemize}[noitemsep]
\item \textbf{Test automation scripts:} Python runners, PowerShell/Bash wrappers for each suite.
\item \textbf{Documented test cases:} CSV files with inputs, expected outputs, and actual results.
\item \textbf{Mutation artifacts:} Unified diff file (\texttt{mutants.diff}) and per-mutant CSV results.
\item \textbf{Test logs:} Per-suite, per-platform execution logs with timestamps.
\item \textbf{Coverage reports:} LLVM-cov and LCOV formatted coverage data.
\item \textbf{Incident reports:} Conformance deviations and portability failures.
\item \textbf{Statistical summaries:} JSON summary files for each enhanced suite.
\item \textbf{Final report:} Comprehensive write-up with figures and analysis.
\end{itemize}

% ══════════════════════════════════════════════════════════════════
\section{Risk Assessment}
% ══════════════════════════════════════════════════════════════════

\begin{table}[H]
\centering
\caption{Testing risks and mitigations.}
\begin{tabularx}{\textwidth}{lXX}
\toprule
\textbf{Risk} & \textbf{Impact} & \textbf{Mitigation} \\
\midrule
MSVC/libFuzzer incompatibility & No fuzzing on Windows & Defer to WSL; document as portability gap \\
Long fuzzing campaigns & Limited time budget per target & Use 60s+180s budgets; prioritize coverage \\
Equivalent mutants inflate scores & Misleading adequacy ratios & Manual equivalence classification for top survivors \\
TunTap Linux-only gating & No official example on Windows & Develop custom loopback harness \\
\bottomrule
\end{tabularx}
\end{table}

\end{document}
